\chapter{Bulk Transfer End to End}\label{ch:transfer}

The bulk-transfer API (\code{fuse/transfer.hpp}, implemented in \code{src/proto/transfer.cpp}) is the most
complete transport and the one that received the full version~8 hardening. This chapter follows a transfer
from configuration to completion.

\section{Interface and shape of a transfer}

The API consists of four blocking calls: \code{send_buffer}, \code{receive_buffer}, \code{send_file} and
\code{receive_file}. Each takes a \code{TransferConfig} (Appendix~\ref{app:config}) and returns a
\code{TransferStatus}. The receiver must be started first, because it has to be bound before the sender's
first message arrives.

The data is split into \code{lanes} equal shards, four by default. Lane $i$ is an independent reliable stream
with these properties:
\begin{itemize}
  \item it uses UDP port \code{base_port}$+i$ on the receiver, so the whole port range must be open;
  \item it has its own thread on the sender;
  \item it has its own reader thread on the receiver;
  \item it has its own sequence space, window, registry, congestion controller and keys.
\end{itemize}
A configuration whose last port would exceed 65,535 is rejected\vviii; before version~8 the sum wrapped around
to ports near zero.

On the receiver, a lane's reader thread does not write data itself. It hands each new block to one of $K$ writer
threads (\code{writer_threads}, two by default), or to a dedicated thread for retransmitted blocks, through
lock-free rings (Figure~\ref{fig:pipeline}). The writers decrypt and write each block directly to its final
position in the output. Only when a writer reports success does the reader admit the block to its window, so
every acknowledgement the sender sees covers bytes that are both authentic and already written.

\begin{figure}[H]
\centering
\fitwidth{\begin{tikzpicture}[font=\small]
  \node[io,minimum height=1.2cm] (sock) at (0,0) {socket\\lane $i$};
  \node[core,minimum height=2.4cm,minimum width=2.4cm] (rd) at (3.4,0) {reader\\(window,\\ACK, NACK)};
  \node[box,minimum width=2.4cm] (w1) at (8,1.5) {writer 1};
  \node[lbl] at (8,0.75) {$\vdots$};
  \node[box,minimum width=2.4cm] (w2) at (8,0) {writer $K$};
  \node[box,minimum width=2.4cm] (rt) at (8,-1.5) {retransmit consumer};
  \node[io,minimum height=2.4cm,minimum width=2.4cm] (sink) at (12.4,0) {output\\(memory or\\\texttt{.part} file)};
  \draw[arr] (sock) -- (rd);
  \draw[arr] (rd.east|-w1) -- node[above,lbl]{new blocks} (w1.west);
  \draw[arr] (rd.east|-w2) -- node[above,lbl]{new blocks} (w2.west);
  \draw[arr] (rd.east|-rt) -- node[above,lbl]{retransmissions} (rt.west);
  \draw[arr] (w1.east) -- (sink.west|-w1); \draw[arr] (w2.east) -- (sink.west|-w2); \draw[arr] (rt.east) -- (sink.west|-rt);
  \node[lbl,below=3mm of rt] {each arrow from the reader is a ring; a matching ring carries each verdict back};
\end{tikzpicture}}
\caption{Receive pipeline of one lane. Solid arrows carry datagrams and dashed arrows carry verdicts back.
Every ring has one producer and one consumer, so no mutex is needed.}\label{fig:pipeline}
\end{figure}

New blocks are routed to writer $\lfloor \mathit{seq}/64\rfloor \bmod K$: whole runs of 64 consecutive blocks go
to the same writer, which can then write each run with a single \code{pwritev}. Round-robin routing
($\mathit{seq}\bmod K$) left every writer's blocks non-adjacent, and was measured to be two to four times slower
on a real disk. A full ring never blocks the reader: the datagram is dropped and recovered by the normal loss
recovery.

\section{Opening a lane}\label{sec:handshake}

Figure~\ref{fig:handshake} shows how a lane is opened. The sender draws a random 64-bit lane nonce from a
secure source and sends a StreamStart. Under a pre-shared key, the StreamStart carries the session salt and the
control trailer.

\begin{figure}[H]
\centering
\fitwidth{\begin{tikzpicture}[font=\small,yscale=0.9]
  \node[box] at (0,0.4) {Sender lane $i$};
  \node[box] at (9,0.4) {Receiver lane $i$};
  \draw[thick,black!40] (0,0) -- (0,-11.6); \draw[thick,black!40] (9,0) -- (9,-11.6);
  \draw[msg] (0,-0.6) -- node[above,sloped]{StreamStart(nonce, token $=0$)} (9,-1.2);
  \node[lbl,right] at (9.1,-1.3) {token wrong:\\commit nothing};
  \draw[msg] (9,-1.9) -- node[above,sloped]{Retry(nonce, token)\vviii} (0,-2.5);
  \draw[msg] (0,-2.9) -- node[above,sloped]{StreamStart(nonce, token)} (9,-3.5);
  \node[lbl,right] at (9.1,-3.9) {verify GMAC\vviii\\check sizes, prepare sink\\claim lane};
  \draw[msg] (9,-4.5) -- node[above,sloped]{ResumeRanges(nonce)[parts]} (0,-5.1);
  \draw[msg] (9,-5.4) -- node[above,sloped]{Ack(nonce)} (0,-6.0);
  \draw[msg] (0,-6.5) -- node[above,sloped]{DATA \ldots} (9,-7.1);
  \draw[msg] (9,-7.4) -- node[above,sloped]{Ack / Nack} (0,-8.0);
  \draw[msg,dashed] (0,-8.5) -- node[above,sloped]{FileDigest (resumed only)} (9,-9.1);
  \draw[msg,dashed] (9,-9.4) -- node[above,sloped]{FileDigest (echo)} (0,-10.0);
  \draw[msg] (0,-10.5) -- node[above,sloped]{StreamClose(finished) $\times3$} (9,-11.1);
\end{tikzpicture}}
\caption{The life of a lane. Every reply carries the lane nonce; under a pre-shared key every message except
Retry carries a control trailer.}\label{fig:handshake}
\end{figure}

\paragraph{Address validation\vviii.} On receiving a StreamStart, the receiver first computes
\[
\mathit{token}=\mathrm{SipHash\text{-}2\text{-}4}_{s}(\mathit{family}\,\|\,\mathit{port}\,\|\,\mathit{address}\,\|\,\mathit{nonce}),
\]
where $s$ is a 128-bit secret held only by this receiver lane. If the StreamStart does not carry this token, the
receiver replies with a 20-byte Retry and does nothing else. It allocates nothing, claims nothing and sends
nothing larger than the request. Only a sender that genuinely receives at its claimed address can echo the
token. A spoofed StreamStart can therefore neither take over the lane nor turn the receiver into an amplifier
aimed at a victim. This is the same mechanism as QUIC's Retry~\cite{rfc9000}.

\paragraph{Key confirmation\vviii.} When a StreamStart with a valid token arrives and the lane has no session
yet, the receiver authenticates it before allocating anything. It derives the session key from its own
pre-shared key and the StreamStart's salt, creates a trial control authenticator, and verifies the trailer. Two
outcomes are possible:
\begin{itemize}
  \item \emph{Verification fails.} This happens when the keys differ, when only one side has a key, or when a
    plaintext receiver meets an encrypting sender. The receiver answers, unauthenticated and at most once per
    100~ms, with StreamClose(reason~2). The sender then reports \code{AuthFailed} immediately, instead of after
    twenty failed data blocks or a timeout.
  \item \emph{Verification succeeds.} The receiver validates the sizes in an order that cannot overflow, namely
    $\mathit{total\_bytes}\le\mathit{file\_total}$ and
    $\mathit{base\_offset}\le\mathit{file\_total}-\mathit{total\_bytes}$. The first lane to succeed then sizes
    the output with an atomic compare-and-swap, and every other lane must agree on the total. If the claimed
    size cannot be allocated, the receiver fails with \code{ResourceLimit}.
\end{itemize}

\paragraph{Answer and retries.} The receiver then records the lane nonce and the sender's address, computes the
byte ranges of the lane it still needs (all of it, unless it is resuming), and sends them as ResumeRanges
followed by Ack messages. The sender considers the handshake complete only once it has an Ack and every
ResumeRanges part carrying its nonce. It retries its StreamStart with exponential backoff from 1~ms to 200~ms,
bounded by \code{connect_timeout_ms} (10~s by default); a retried StreamStart with the right nonce makes the
receiver re-send its answer. The ranges are checked for order, overlap and bounds before use, and an
inconsistent answer makes the sender send the whole lane instead.

\section{The sender's main loop}\label{sec:sendloop}

After the handshake the sender lane runs one loop until every byte is acknowledged or the lane gives up. Each
iteration performs five steps.

\begin{enumerate}
  \item \textbf{Stop checks.} The loop first handles user cancellation and the shared \code{stop_all} flag, which
    is set when a sibling lane has heard the peer abort.
  \item \textbf{Window computation.} It advances the congestion controller and computes the bytes it may have in
    flight:
    \[
    \mathit{window\_bytes}=\min\bigl(\max(4,\mathit{cwnd})\cdot\mathit{base\_block},\;\max(4,\mathit{rwnd})\cdot\mathit{block}\bigr).
    \]
    Here $\mathit{cwnd}$ is counted in units of the \emph{configured} block size\vviii{} and $\mathit{rwnd}$ is
    the receiver's advertised window\vviii. The floor of four keeps a zero window from wedging the flow.
  \item \textbf{Sending a batch.} If there is room in the window and in the 1,024-slot registry, and the pacing
    time has passed, the sender builds a batch:
    \begin{itemize}
      \item It stores each block in the registry, seals it, encodes the DATA datagram into a staging buffer, and
        sends the batch with segmentation offload, falling back to individual sends.
      \item The batch is limited by the offload budget (about 60~KB), the registry, the window and a burst cap.
        The cap is 16 datagrams once the path has a real round-trip time (above 1~ms), because a 64-datagram
        line-rate burst overflows shallow router buffers.
      \item The next batch is paced (Section~\ref{sec:cc}).
    \end{itemize}
  \item \textbf{Reading control traffic.} It drains up to 64 control datagrams. Under a pre-shared key, any
    datagram without a valid trailer is discarded before it is even parsed. Otherwise:
    \begin{itemize}
      \item an \emph{Ack} with the right nonce is applied with \code{apply_ack}, its \code{rwnd} recorded, and a
        new echoed send time fed to the RTT estimator;
      \item a \emph{Nack} marks a loss in the controller and resends each listed block;
      \item a \emph{StreamClose(aborted)} ends the lane with \code{PeerAborted};
      \item a \emph{PathChallenge} is answered, and a \emph{PathResponse} completes a pending path validation
        (Section~\ref{sec:path}).
    \end{itemize}
  \item \textbf{Timeout and adaptation.} It runs the retransmission timeout and the block-size adaptation
    (Sections~\ref{sec:recovery} and~\ref{sec:blocksize}), and gives up with \code{Timeout} if no acknowledgement
    has confirmed new data for \code{timeout_ms}.
\end{enumerate}

\paragraph{RTT samples from repeated echoes.} A repeated echo is never treated as a new RTT sample. The receiver
only updates its echoed send time when a new highest block arrives, so its periodic acknowledgements repeat the
same value. Counting each repeat as a fresh sample manufactured an ever-growing fake RTT whenever the sender was
stalled on a gap, which stretched the timer exactly when it needed to stay responsive.

\paragraph{The empty case.} A lane with nothing to send still needs a completion signal: an empty send, an
uneven split, or a resumed lane that is already complete. The lane sends a single empty block marked LastBlock.

\section{The life of one packet}\label{sec:life}

To make the pieces concrete, we follow one encrypted block, number 4,711 of lane~2, from the sender's memory to
the receiver's disk.

\begin{enumerate}
  \item \textbf{Admission.} The lane's congestion and receive windows have room and the pacing time has passed.
    The block's bytes are copied into registry slot $4711 \bmod 1024$, together with their offset, length and
    flags.
  \item \textbf{Key and nonce.} The key schedule selects the cipher for epoch $4711\gg22=0$. The nonce is
    $\langle 2, 4711, \text{0x0000}\rangle$, and the AAD is $\langle 2, 4711, \mathit{offset}, \mathit{flags}
    \wedge\lnot\text{0x01}\rangle$.
  \item \textbf{Sealing.} AES-256-GCM encrypts the payload into a scratch buffer and appends the tag; if sealing
    fails, the batch stops.
  \item \textbf{Framing.} The datagram is the outer header, the 21-byte block header with
    \code{payload_len} including the tag, the send time plus the lane's secret offset, and the ciphertext.
  \item \textbf{Transmission.} It joins up to 15 other datagrams in the staging buffer, and the batch leaves in
    one \code{sendmsg} with \code{UDP_SEGMENT}. The kernel splits the buffer into individual datagrams.
  \item \textbf{Reception.} On the receiver, the lane's reader collects it with up to 63 others in one
    \code{recvmmsg}, peeks at the type, and decodes the header. It discards the block if the block is below the
    base, beyond the window, already received, or currently held by a writer. Otherwise it copies the datagram
    into the ring of writer $\lfloor4711/64\rfloor\bmod K$.
  \item \textbf{Decryption and write.} The writer selects its own cipher for the lane and epoch, rebuilds the
    AAD, and opens the block. On success it checks that
    $\mathit{offset}\le\mathit{lane\_total}$ and $\mathit{len}\le\mathit{lane\_total}-\mathit{offset}$, merges
    the block with any adjacent blocks in its batch, and writes the run with one \code{pwritev} at
    $\mathit{lane\_base}+\mathit{offset}$. It posts the verdict to the reader's completion ring.
  \item \textbf{Admission to the window.} The reader collects the verdict and calls \code{on_receive}. The bit
    for 4,711 is set; if it was the first gap, the base slides forward. If the block came from a new source
    address, a path challenge is sent there (Section~\ref{sec:path}).
  \item \textbf{Acknowledgement.} After eight newly delivered blocks, or 200~\textmu s, or a window update, the
    reader sends an Ack. The Ack carries the new base, the bitmask, the echoed send time, the nonce and the
    current receive window, and is signed with the control trailer.
  \item \textbf{Release.} The sender verifies the trailer and applies the Ack, releasing slot $4711\bmod 1024$.
    The block's bytes now count towards progress, and the window may grow.
\end{enumerate}

\section{Loss recovery}\label{sec:recovery}

Fuse recovers losses through three cooperating mechanisms.

\paragraph{Negative acknowledgements.} The receiver reports gaps that have persisted beyond a reorder delay and
were not reported within a re-NACK interval. Both intervals scale with the receiver's own estimate of the RTT:
\begin{itemize}
  \item the reorder delay is half the RTT, clamped to 1--100~ms;
  \item the re-NACK interval is 1.5 times the RTT, clamped to 20--500~ms;
  \item before any estimate exists, conservative wide-area defaults of 3~ms and 40~ms apply.
\end{itemize}
The receiver measures the RTT purely on its own clock, as the delay between sending a NACK and receiving the
retransmission it requested, so no comparison between the two hosts' clocks is needed. Scaling matters: with
fixed loopback-tight values, a wide-area path received a duplicate NACK every few milliseconds while the
retransmission was still 50~ms away.

\paragraph{Selective acknowledgement\vviii.} Every Ack confirms blocks above the base through its bitmask, so a
NACK or timeout never resends a block the receiver already holds.

\paragraph{Retransmission timeout\vviii.} If the base has not advanced for one RTO (Section~\ref{sec:cc}), the
sender resends the base and up to seven further \emph{known holes}: unacknowledged blocks below
\code{received_end}, the highest block the receiver is known to hold. After a burst loss, the rest of the burst
is therefore repaired in one timeout rather than one timeout per block, without resending blocks that are merely
in flight. A send that the kernel refuses, for example because the socket buffer is full, does not count as
progress; otherwise the block would never be retransmitted and the transfer would stall.

\section{Flow control}\label{sec:flow}

Before version~8 there was no flow control. A receiver whose writers fell behind a fast sender (because of a
slow disk, say) silently dropped blocks at its full rings. The sender interpreted the losses as network
congestion: NACKs, retransmissions and a halved window.

Now each Ack advertises the receive window, which is the free space in the lane's writer rings plus the blocks
already handed over but not yet written, capped at 1,024. The sender keeps its bytes in flight within it. When
the advertised window was low (below a quarter of the maximum) and has since at least doubled, the receiver
sends an Ack immediately as a \emph{window update}. Without the update, a sender that paused for the receiver
would sit idle until the next periodic acknowledgement. This is the behaviour that removed a throughput
regression found while the feature was being developed.

\section{Adaptive block size and the MTU}\label{sec:blocksize}

The sender starts with the configured block size, 1,200 bytes by default, and doubles it after eight batches
in a row that were confirmed by acknowledgements, up to a \emph{growth ceiling}\vviii. Any loss returns the
block size to the configured value. Growth is counted in acknowledged batches rather than loop iterations, which
can accumulate many times within one round trip, before any acknowledgement could have arrived.

The growth ceiling is computed from the kernel's route MTU:
\[
\mathit{ceiling}=\min\bigl(\mathit{block\_ceiling},\;\mathit{mtu}-\mathit{ip\_udp}-31-\mathit{tag}\bigr),
\]
and is never below the starting size. An unknown MTU means no growth at all. On loopback, which reports a
64~KiB MTU, blocks may still grow to 16~KiB; on an ordinary 1,500-byte path they stop at about 1,450 bytes.
Before version~8, blocks grew to 16~KiB on any path. Above the MTU every such datagram is fragmented, and losing
one fragment loses the whole datagram.

On top of the MTU bound, a loss that occurs soon after a growth step makes the previous size the new ceiling,
so on a path that drops large datagrams the size settles instead of oscillating. In the spirit of
DPLPMTUD~\cite{rfc8899}, the larger size is probed again after a clean stretch, and that stretch doubles after
each failed probe, up to 4,096 acknowledged batches. One transient loss therefore does not cap the lane
permanently.

The congestion window counts units of the \emph{configured} block size. Its cap is the number of units that
1,024 blocks of the largest permitted size would occupy, so a long fat path can still be filled, but only by
earning the window through slow start and additive increase. Previously the window counted blocks of whatever
size they had grown to, so the same window could hold up to 13 times as many bytes. Slow start starts from 128
units, about the 150~KB a lane used to begin with.

\section{Path validation}\label{sec:path}

A receiver whose peer is behind a NAT may see the sender's source address change in the middle of a transfer
when the NAT rebinds. Before version~8, the bulk transfer simply re-anchored its peer address to the source of
any accepted block. An attacker could then redirect a lane's acknowledgements, or reflect traffic at a third
party, with a single spoofed datagram.

Now, when authentic session traffic arrives from a new source, the side that sees it sends a PathChallenge with
a fresh random 64-bit value to that source. It switches its peer address only when a matching PathResponse comes
back \emph{from that source}. One candidate is validated at a time, and challenges are resent at most every
50~ms. Data from the new source is accepted in the meantime, so a genuine rebinding costs no throughput. The
mechanism is symmetric: the sender validates a new source of Acks in the same way.

\section{Finishing a lane}

The receiver knows a lane is complete when it has seen the block carrying LastBlock and the base has moved past
it. The flag is trusted only after the block has authenticated.

The receiver does not exit at once. It keeps sending acknowledgements on its normal schedule in case its last
ones were lost, and it leaves at the first of three events:
\begin{itemize}
  \item the sender's StreamClose(finished) arrives;
  \item the sender has been silent for 400~ms;
  \item 1.5~s have passed since completion.
\end{itemize}
It then sends eight final acknowledgements. The lane's reported duration ends at completion rather than after
this linger, so throughput figures are not diluted.

For a resumed transfer, lane~0 also waits for the sender's whole-file digest. The sender computes the digest in a
background thread that starts right after the handshake, overlapped with sending the missing data. It re-sends
the digest every 50~ms until the receiver echoes it, and then sends StreamClose(finished) three times.

\section{State machines of the two lane roles}

Figures~\ref{fig:sendfsm} and~\ref{fig:recvfsm} summarise the behaviour described above.

\begin{figure}[H]
\centering
\fitwidth{\begin{tikzpicture}[>={Stealth[length=2.2mm]},font=\small]
  \node[st,minimum width=3.2cm] (hs) at (0,0) {Handshake\\(StreamStart, backoff)};
  \node[st,minimum width=3.2cm] (sending) at (6.5,0) {Sending\\(window, pacing, RTO)};
  \node[st,minimum width=2.4cm,fill=green!10] (ok) at (12.4,0) {Done: Ok};
  \node[st,minimum width=3.2cm] (digest) at (6.5,-3.4) {Digest\\(lane 0, resumed)};
  \node[st,minimum width=2.4cm,fill=red!8] (fail) at (0,-3.4) {Failed};
  \draw[->,thick] (-2.6,0) -- (hs);
  \draw[->,thick] (hs) -- node[lbl,below]{Ack and all\\ResumeRanges parts} (sending);
  \draw[->,thick] (hs.north west) to[out=100,in=80,looseness=1.6] node[lbl,above]{Retry: echo the token} (hs.north east);
  \draw[->,thick] (hs) -- node[lbl,left,align=right]{auth-failed close,\\version mismatch,\\connect timeout,\\cancel} (fail);
  \draw[->,thick] (sending) -- node[lbl,above]{all acknowledged\\(fresh)} (ok);
  \draw[->,thick] (sending) -- node[lbl,right]{all acknowledged (resumed)} (digest);
  \draw[->,thick] (digest) -| node[lbl,below,pos=0.25]{digest echoed} (ok);
  \draw[->,thick] (sending) -- node[lbl,below,sloped,pos=0.45]{abort, timeout, cancel} (fail);
  \draw[->,thick] (digest) -- node[lbl,below]{abort, timeout} (fail);
\end{tikzpicture}}
\caption{Sender lane. Every transition to \emph{Failed} except a received abort sends StreamClose(aborted) to the peer.}\label{fig:sendfsm}
\end{figure}

\begin{figure}[H]
\centering
\fitwidth{\begin{tikzpicture}[>={Stealth[length=2.2mm]},font=\small]
  \node[st,minimum width=2.8cm] (wait) at (0,0) {Waiting\\(no session)};
  \node[st,minimum width=3.2cm] (recv) at (6,0) {Receiving\\(window, ACK, NACK)};
  \node[st,minimum width=2.8cm] (linger) at (12,0) {Complete,\\lingering};
  \node[st,minimum width=2.8cm,fill=green!10] (ok) at (12,-3.4) {Done: Ok\\or Incomplete};
  \node[st,minimum width=2.8cm,fill=red!8] (fail) at (0,-3.4) {Failed};
  \draw[->,thick] (-2.4,0) -- (wait);
  \draw[->,thick] (wait.north west) to[out=100,in=80,looseness=1.6] node[lbl,above]{no token: Retry; bad auth: close(2)} (wait.north east);
  \draw[->,thick] (wait) -- node[lbl,below]{valid token,\\authentic, sizes valid} (recv);
  \draw[->,thick] (recv) -- node[lbl,above]{base passes\\the last block} (linger);
  \draw[->,thick] (linger) -- node[lbl,right,align=left]{close(finished), or\\400 ms quiet, or 1.5 s} (ok);
  \draw[->,thick] (wait) -- node[lbl,left,align=right]{unallocatable size:\\ResourceLimit} (fail);
  \draw[->,thick] (recv) -- node[lbl,below,sloped,pos=0.45]{abort, idle timeout, cancel, I/O error} (fail);
\end{tikzpicture}}
\caption{Receiver lane. The final status is classified as \texttt{AuthFailed} if most blocks failed authentication.}\label{fig:recvfsm}
\end{figure}

\section{Cancellation, abort, timeout and combining statuses}

\paragraph{Cancellation.} \code{TransferConfig::cancel} is an atomic flag that every lane polls. Its loop
cadence is bounded by short socket timeouts, at most 20~ms on the receiver. When the flag is set, a lane sends
StreamClose(aborted) three times and returns \code{Cancelled}, and its siblings stop as well. The peer returns
\code{PeerAborted} instead of waiting out its timeout.

\paragraph{Aborts before a session exists.} Under plaintext, an abort is accepted with any nonce, so a sender
cancelled before its StreamStart arrived can still say ``I am not coming''. Under a pre-shared key, such an abort
cannot be verified and is ignored.

\paragraph{Timeouts.} The receiver gives up when no datagram has arrived for \code{timeout_ms} plus one second.
The sender gives up when no acknowledgement has confirmed new data for \code{timeout_ms}.

\paragraph{Detecting an active attack.} If at least 20 blocks have been seen and more than half of them failed
authentication, the receiver reports \code{AuthFailed}. An active tampering attack fails most blocks, which is
distinguishable from ordinary loss, and the caller should learn of it as such.

\paragraph{Combining lane statuses.} The call's result is the most severe lane status, so that the caller sees
the most actionable problem rather than whichever lane finished last. The ranking, from least to most severe,
is:
\begin{quote}\small
Ok $<$ Incomplete $<$ Timeout $<$ PeerAborted $<$ Cancelled $<$ AuthFailed $<$ VerifyFailed $<$ ResourceLimit
$<$ IoError $<$ SocketError $<$ VersionMismatch $<$ ConfigError, Unsupported.
\end{quote}

\section{Resumable file reception}\label{sec:resume}

\code{receive_file} writes into \code{<path>.part} and renames it over \code{<path>} only after every byte has
arrived and been synchronised to disk, so the destination file is always either untouched or complete. Behind
the data, the \code{.part} file carries a trailer
\[
[\,\text{data}\,][\,\text{bitmap: one bit per chunk}\,][\,\mathit{file\_id},\mathit{total},\mathit{chunk},\mathit{nchunks},\mathit{checksum},\text{``FUSERSM1''}\,].
\]
Chunks are 64~KiB, doubled for very large files until there are at most $2^{23}$ of them. They are used rather
than sequence numbers because the mapping from sequence number to offset changes with the block size, both
during and between sessions. The trailer is rewritten every \code{resume_checkpoint_ms} (one second) and
whenever a transfer ends without completing.

A later \code{receive_file} of the same file adopts the partial copy. Its StreamStart must carry the same
\code{file_id} and size; the receiver then asks only for the missing chunks, merging the closest ranges if the
list would exceed 4,096. At the end it verifies the whole file against the sender's digest. On a mismatch it
reports \code{VerifyFailed} and \emph{discards} the copy, since another resume would fail the same way.

The file identity hashes the file's name, size, modification time, inode change time, inode number and
device\vviii. Before version~8 it used only name, size and modification time, which missed a file rewritten in
place with its modification time set back (by \code{touch -r} or a restore tool); a resume would then have
stitched stale bytes into the result. The inode change time cannot be set back by user code.

\section{Timers}

\begin{table}[H]
\centering\small
\caption{Every timer of the bulk transfer.}\label{tab:timers}
\begin{tabularx}{\linewidth}{>{\raggedright\arraybackslash}p{4.4cm}>{\raggedright\arraybackslash}p{3.6cm}X}
\toprule
\textbf{Timer} & \textbf{Value} & \textbf{Purpose} \\
\midrule
StreamStart retry & 1~ms doubling to 200~ms & handshake on any path RTT \\
Connect timeout & \code{connect_timeout_ms} (10~s) & fail fast if no receiver \\
No-progress timeout & \code{timeout_ms} (120~s) & sender: no acknowledged progress \\
Receiver idle limit & \code{timeout_ms} $+$ 1~s & receiver: no datagrams at all \\
Periodic ACK & every 8 blocks or 200~\textmu s & feedback cadence \\
Quiet ACK limit & 3 to a silent peer\vviii & no endless reflection \\
Reorder delay & RTT/2, 1--100~ms (3~ms initially) & late is not lost \\
Re-NACK interval & 1.5 RTT, 20--500~ms (40~ms initially) & no duplicate repair \\
RTO & Section~\ref{sec:cc}, 5~ms--1~s, backoff $\le 2^6$ & tail and burst losses \\
Path challenge resend & 50~ms\vviii & validate a new address \\
Version / auth reply limit & 100~ms\vviii & no reflection amplifier \\
Digest resend & 50~ms & resumed transfers \\
Completion linger & close, 400~ms quiet, or 1.5~s & lost final ACKs \\
Resume checkpoint & \code{resume_checkpoint_ms} (1~s) & crash-safe progress \\
\bottomrule
\end{tabularx}
\end{table}
